% !TeX root = ../main.tex
% !TeX encoding = utf8
\section{NỘI DUNG MỞ RỘNG: PERCENTILE CI VÀ TÍNH BẤT BIẾN}

\subsection{Khoảng tin cậy Bootstrap phân vị (Percentile CI)}

Bên cạnh Normal CI và Basic CI, Percentile CI lấy trực tiếp các phân vị của phân phối Bootstrap làm hai cận khoảng \cite{efron1993}.

\subsubsection{Cơ sở và công thức}
Percentile CI đặt giả định rằng phân phối của bản sao Bootstrap $\hat{\theta}^*$ xấp xỉ trực tiếp phân phối lấy mẫu của ước lượng điểm $\hat{\theta}_n$ \cite{efron1993}. Tức là, các phân vị thực nghiệm của tập hợp $\{\hat{\theta}^{*(1)}, \dots, \hat{\theta}^{*(B)}\}$ được sử dụng trực tiếp làm cận của khoảng tin cậy mà không qua phép biến đổi độ lệch nào \cite{efron1993}:

\begin{equation}
	CI_{\text{Percentile}} = \left[ q_{\alpha/2}^*, \quad q_{1-\alpha/2}^* \right]
\end{equation}
trong đó $q_{\alpha/2}^*$ và $q_{1-\alpha/2}^*$ tương ứng là phân vị thực nghiệm mức $\alpha/2$ và $1-\alpha/2$ của các bản sao Bootstrap \cite{buithuytrang2026}.
\subsection{Khoảng tin cậy Bootstrap BCa}

Khoảng tin cậy Bootstrap BCa (Bias-Corrected and Accelerated) điều chỉnh
các phân vị Bootstrap để phản ánh độ chệch và độ bất đối xứng của phân phối
ước lượng \cite{efron1993, davison1997}.

Gọi $\hat{\theta}_n$ là ước lượng từ mẫu gốc và
$\hat{\theta}^{*(1)},\ldots,\hat{\theta}^{*(B)}$ là các bản sao Bootstrap.
Hệ số hiệu chỉnh độ chệch được ước lượng bởi:
\begin{equation}
	\widehat{z}_0
	=
	\Phi^{-1}\left(
	\frac{1}{B}\sum_{b=1}^{B}
	\mathbf{1}\left\{
	\hat{\theta}^{*(b)} < \hat{\theta}_n
	\right\}
	\right),
\end{equation}
trong đó $\Phi^{-1}$ là hàm phân vị của phân phối Chuẩn tắc.

Tiếp theo, gọi $\hat{\theta}_{(i)}$ là giá trị thống kê tính khi bỏ quan sát
thứ $i$ khỏi mẫu, và
$\bar{\theta}_{(.)}=n^{-1}\sum_{i=1}^{n}\hat{\theta}_{(i)}$. Hệ số gia tốc
được ước lượng bằng:
\begin{equation}
	\widehat{a}
	=
	\frac{
		\displaystyle\sum_{i=1}^{n}
		\left(\bar{\theta}_{(.)}-\hat{\theta}_{(i)}\right)^3
	}{
		\displaystyle
		6\left[
		\sum_{i=1}^{n}
		\left(\bar{\theta}_{(.)}-\hat{\theta}_{(i)}\right)^2
		\right]^{3/2}
	}.
\end{equation}

Với mức ý nghĩa $\alpha$, hai mức phân vị điều chỉnh được xác định bởi:
\begin{equation}
	p_L
	=
	\Phi\left[
	\widehat{z}_0+
	\frac{\widehat{z}_0+z_{\alpha/2}}
	{1-\widehat{a}(\widehat{z}_0+z_{\alpha/2})}
	\right],
\end{equation}
\begin{equation}
	p_U
	=
	\Phi\left[
	\widehat{z}_0+
	\frac{\widehat{z}_0+z_{1-\alpha/2}}
	{1-\widehat{a}(\widehat{z}_0+z_{1-\alpha/2})}
	\right],
\end{equation}
trong đó $z_p=\Phi^{-1}(p)$. Hai mức $p_L$ và $p_U$ nói chung không có tổng bằng $1$. Khoảng tin cậy BCa được tính bằng:
\begin{equation}
	CI_{\mathrm{BCa}}
	=
	\left[
	q_{p_L}^{*},
	q_{p_U}^{*}
	\right].
\end{equation}

Các công thức trên yêu cầu $\widehat z_0$ hữu hạn, mẫu Jackknife không suy biến và các mẫu số của phép hiệu chỉnh khác $0$. Nếu các điều kiện này không thỏa mãn, cần kiểm tra tính khả dụng của BCa trước khi báo cáo khoảng.

Percentile CI lấy trực tiếp hai phân vị Bootstrap thông thường, còn BCa
điều chỉnh các mức phân vị để xét đến độ chệch và độ bất đối xứng.
BCa có thể cải thiện độ chính xác trong những điều kiện thích hợp, nhưng
không bảo đảm tốt hơn trong mọi trường hợp; kết quả còn phụ thuộc vào dữ
liệu, thống kê đang xét và số lần lấy mẫu Bootstrap
\cite{efron1993, davison1997}.
\subsection{So sánh Normal CI, Basic CI và Percentile CI}

Bảng \ref{tab:trio_comparison} so sánh ba phương pháp khoảng tin cậy Bootstrap \cite{efron1993, davison1997}. Đặt $q_L^*=q_{\alpha/2}^*$ và $q_U^*=q_{1-\alpha/2}^*$.

\begin{table}[htbp]
	\centering
	\renewcommand{\arraystretch}{1.4}
	\setlength{\tabcolsep}{5pt}
	\caption{So sánh Normal CI, Basic CI và Percentile CI}
	\label{tab:trio_comparison}
	\small
	\setlength{\tabcolsep}{4pt}
	\begin{tabularx}{\textwidth}{>{\raggedright\arraybackslash}p{2.5cm}XXX}
		\toprule
		\textbf{Tiêu chí} & \textbf{Normal CI} & \textbf{Basic CI} & \textbf{Percentile CI} \\
		\midrule
		\textbf{Công thức} & $\hat{\theta}_n \pm z_{1-\alpha/2}\widehat{se}^*$ \cite{buithuytrang2026} & $[2\hat{\theta}_n-q_U^*,\;2\hat{\theta}_n-q_L^*]$ \cite{buithuytrang2026} & $[q_L^*,q_U^*]$ \cite{buithuytrang2026} \\
		\textbf{Cơ sở} & Dựa trên xấp xỉ chuẩn của ước lượng \cite{efron1993}. & Dựa trên xấp xỉ phân phối sai số Bootstrap \cite{davison1997}. & Dùng trực tiếp các phân vị Bootstrap \cite{efron1993}. \\
		\textbf{Phân phối lệch phải} & Khoảng đối xứng quanh $\hat{\theta}_n$ \cite{efron1993}. & Quantile Flipping có thể tạo khoảng bất đối xứng \cite{davison1997}. & Hai cận theo hướng phân phối Bootstrap \cite{efron1993}. \\
		\textbf{Bất biến qua $g(\theta)$} & Nói chung không bất biến \cite{efron1993}. & Nói chung không bất biến \cite{davison1997}. & Bất biến qua biến đổi đơn điệu, nếu quy tắc phân vị tương thích \cite{efron1993}. \\
		\textbf{Độ chính xác tiệm cận} & Bậc nhất, tùy điều kiện \cite{efron1993}. & Bậc nhất, tùy điều kiện \cite{davison1997}. & Bậc nhất, tùy điều kiện \cite{efron1993}. \\
		\bottomrule
	\end{tabularx}
\end{table}

\subsection{Tính bất biến dưới phép biến đổi tham số $g(\theta)$}

\begin{definition}[\textbf{Tính Bất biến của Khoảng Tin cậy}]
	Một phương pháp xây dựng khoảng tin cậy cho tham số $\theta$ được gọi là \textbf{bất biến (invariant)} dưới phép biến đổi đơn điệu $g(\cdot)$ nếu khoảng tin cậy cho tham số mới $\eta = g(\theta)$ đúng bằng việc áp dụng ánh xạ $g(\cdot)$ lên các cận của khoảng tin cậy cho $\theta$ \cite{efron1993}:
	\begin{equation}
		CI_{g(\theta)} =
    \begin{cases}
      [g(L_{\theta}),g(U_{\theta})], & g\ \text{đồng biến},\\
      [g(U_{\theta}),g(L_{\theta})], & g\ \text{nghịch biến}.
    \end{cases}
	\end{equation}
	\end{definition}
		
		\subsubsection{Tính bất biến của Percentile CI}
		Giả sử $g(\cdot)$ là một hàm số đồng biến nghiêm ngặt (strictly monotonic increasing function), chẳng hạn như phép biến đổi Logarithm $g(x) = \ln(x)$ hoặc phép biến đổi Lũy thừa $g(x) = x^2$ với $x > 0$.
		Trong chứng minh này, dùng phân vị thực nghiệm type~1 (nghịch đảo tổng quát của ECDF), không nội suy giữa hai giá trị. Với nội suy tuyến tính type~7, phép biến đổi phi tuyến nói chung không giao hoán với nội suy. Vì $g(\cdot)$ đồng biến, thứ tự các phần tử trong tập hợp bản sao Bootstrap được giữ nguyên hoàn toàn \cite{efron1993}:
		\begin{equation}
			\hat{\theta}^{*(1)} \le \hat{\theta}^{*(2)} \le \dots \le \hat{\theta}^{*(B)} \iff g(\hat{\theta}^{*(1)}) \le g(\hat{\theta}^{*(2)}) \le \dots \le g(\hat{\theta}^{*(B)})
		\end{equation}
		Do đó, phân vị thực nghiệm mức $p$ của tham số mới $g(\hat{\theta}^*)$ chính bằng giá trị $g(\cdot)$ áp dụng lên phân vị thực nghiệm mức $p$ của $\hat{\theta}^*$ \cite{efron1993}:
		\begin{equation}
			q_p^*\left(g(\hat{\theta}^*)\right) = g\left(q_p^*(\hat{\theta}^*)\right)
		\end{equation}
		
		Từ đó, khoảng tin cậy Percentile CI cho tham số $g(\theta)$ là \cite{efron1993}:
		\begin{equation}
			CI_{\text{Percentile}}(g(\theta)) = \left[ q_{\alpha/2}^*(g(\hat{\theta}^*)), \quad q_{1-\alpha/2}^*(g(\hat{\theta}^*)) \right] = \left[ g(q_{\alpha/2}^*), \quad g(q_{1-\alpha/2}^*) \right] \quad (\text{đpcm})
		\end{equation}
		
		\subsubsection{Normal CI và Basic CI dưới phép biến đổi phi tuyến}
		Với phép biến đổi phi tuyến, Normal CI và Basic CI nói chung không bất biến \cite{efron1993, davison1997}:
		\begin{itemize}
			\item \textbf{Normal CI:} Sai số chuẩn và phép biến đổi phi tuyến nói chung không giao hoán; khoảng tính trực tiếp trên thang mới có thể khác với khoảng cũ sau biến đổi \cite{efron1993}.
			\item \textbf{Basic CI:} Phép phản chiếu phân vị quanh ước lượng điểm nói chung không giao hoán với phép biến đổi phi tuyến \cite{davison1997}.
		\end{itemize}
		
		\subsubsection{Kiểm chứng bằng phép biến đổi log}

Dùng cùng các bản sao Bootstrap của mẫu dương $(2,3,4,5,9)$ để so sánh
khoảng tính trên thang $\theta$ rồi lấy log với khoảng tính trực tiếp
trên thang $\eta=\log(\theta)$. Phân vị dùng type~1 trong kiểm chứng này,
để phép biến đổi đơn điệu tăng bảo toàn thứ tự và phân vị thực nghiệm. Đoạn mã là trích đoạn của \texttt{r/Nhom\_D\_TKTT.R}; để tái lập các số ở Bảng~\ref{tab:invariance-log}, chạy toàn bộ tệp với seed và chuỗi mô phỏng đã nêu ở Phần 6.

\begin{lstlisting}[style=Rstyle, caption={Bất biến dưới log}]
invariance_sample <- c(2, 3, 4, 5, 9)
invariance_n <- length(invariance_sample)
invariance_indices <- matrix(
  sample.int(invariance_n, invariance_n * B, replace = TRUE),
  nrow = invariance_n)
bootstrap_theta <- colMeans(matrix(
  invariance_sample[invariance_indices], nrow = invariance_n))
theta_hat <- mean(invariance_sample)
eta_hat <- log(theta_hat)
bootstrap_eta <- log(bootstrap_theta)
q_theta <- quantile(bootstrap_theta, tail_probs, type = 1)
q_eta <- quantile(bootstrap_eta, tail_probs, type = 1)
ci_theta <- c(2 * theta_hat - q_theta[2],
              2 * theta_hat - q_theta[1])
ci_eta_transformed <- log(ci_theta)
ci_eta_direct <- c(2 * eta_hat - q_eta[2],
                   2 * eta_hat - q_eta[1])
max_abs_difference <- max(abs(ci_eta_transformed - ci_eta_direct))
\end{lstlisting}

\begin{table}[htbp]
\centering
\small
\caption{Sai khác lớn nhất giữa hai cách tính khoảng trên thang log}
\label{tab:invariance-log}
\begin{tabular}{lrr}
\toprule
Phương pháp & Sai khác lớn nhất & Khớp trong sai số $10^{-12}$ \\
\midrule
Normal & 0.1566453 & Không \\
Basic & 0.3177451 & Không \\
Percentile & 0 & Có \\
\bottomrule
\end{tabular}
\end{table}

\begin{figure}[htbp]
\centering
\includegraphics[width=0.78\textwidth]{invariance-log-comparison.pdf}
\caption{Sai khác giữa hai cách tính khoảng sau phép biến đổi $\eta=\log(\theta)$.}
\label{fig:invariance-log}
\end{figure}

Trong kiểm chứng này, hai cách tính Percentile trùng nhau đến sai số máy;
Normal và Basic cho kết quả khác nhau sau phép biến đổi log. Các số liệu
chỉ minh họa ba công thức trên mẫu đã chọn. Với biến đổi đơn điệu giảm,
thứ tự hai cận sau biến đổi được đảo lại; kết quả cũng phụ thuộc vào quy
ước tính phân vị của phân phối thực nghiệm rời rạc. Cụ thể, với $g$ giảm, phân vị type~1 tại mức $p$ của mẫu đã biến đổi lấy phần tử gốc có hạng $B+1-\lceil Bp\rceil$, có thể khác hạng $\lceil B(1-p)\rceil$. Vì vậy, phép đảo hai cận chỉ khớp chính xác với khoảng tính lại khi quy tắc chọn hạng tương thích; không suy ra bất biến tuyệt đối cho mọi quy ước phân vị rời rạc.

% ==================== PHẦN 9 ====================
